%-------------------------------------------------------------------------------
%	ТӨЛӨВЛӨГӨӨ, ХИЙГДСЭНИЙ ХУГАЦАА, ЗӨВШӨӨРӨЛ, ШҮҮМЖЛЭГИЙН ХУУДАС
%   «diplom ajluud»-ын Б.Болоржингийн ажилтай ижил бүтэц. Огноо: 2026-2027 намрын хамгаалалтын хуваарь (2026-09-15).
%   Гарын үсэгтэй scan хуудас байвал signatures/huudas1..3.pdf болгож доорх \includepdf-ийн %-ийг арилгана.
%-------------------------------------------------------------------------------

\begin{titlepage}
\small
\noindent Батлав. Удирдагч: \makebox[6cm]{\dotfill} /\supname/
\begin{center}
\vspace*{0.5cm}
\textbf{{\large ТӨГСӨЛТИЙН АЖЛЫН\\ТӨЛӨВЛӨГӨӨ}}\\[0.3cm]
\textbf{Сэдэв: «\ttitle»}\\[0.5cm]
\end{center}
\begin{tabularx}{\textwidth}{|>{\centering\arraybackslash}p{0.6cm}|X|>{\centering\arraybackslash}p{1.6cm}|>{\centering\arraybackslash}p{2.4cm}|}
\hline
\textbf{№} & \textbf{Ажлын бүлгийн нэр} & \textbf{Эзлэх хувь} & \textbf{Дуусах хугацаа} \\ \hline
1 & Оршил & 5\% & 2026-09-30 \\ \hline
2 & Холбоотой ажлуудын судалгаа & 12\% & 2026-09-30 \\ \hline
3 & Арга зүй ба дэвшүүлсэн санаа & 12\% & 2026-10-28 \\ \hline
4 & Өгөгдөл боловсруулалт & 10\% & 2026-10-28 \\ \hline
5 & Загвар хөгжүүлэлт ба урсгал & 15\% & 2026-10-28 \\ \hline
6 & Туршилтын системийн хэрэгжүүлэлт & 12\% & 2026-11-25 \\ \hline
7 & Үр дүн ба хэлэлцүүлэг & 16\% & 2026-11-25 \\ \hline
8 & Ёс зүйн асуудал & 8\% & 2026-11-25 \\ \hline
9 & Дүгнэлт & 10\% & 2026-12-16 \\ \hline
\end{tabularx}

\vspace{0.6cm}
\noindent Төлөвлөгөөг боловсруулсан оюутан: \makebox[4cm]{\dotfill} /\shortname/
\end{titlepage}

\begin{titlepage}
\small
\begin{center}
\textbf{{\large ТӨГСӨЛТИЙН АЖЛЫН ХИЙГДСЭНИЙ\\ХУГАЦАА}}\\[0.5cm]
\end{center}
\begin{tabularx}{\textwidth}{|>{\centering\arraybackslash}p{0.6cm}|X|>{\centering\arraybackslash}p{2.4cm}|>{\centering\arraybackslash}p{3.0cm}|}
\hline
\textbf{№} & \textbf{Хийж гүйцэтгэх ажил} & \textbf{Биелэх хугацаа} & \textbf{Удирдагчийн гарын үсэг} \\ \hline
1 & Үзлэг 1 & 2026-09-30 & \\ \hline
2 & Үзлэг 2 & 2026-10-28 & \\ \hline
3 & Үзлэг 3 & 2026-11-25 & \\ \hline
4 & Урьдчилан хамгаалалт & 2026-12-23 & \\ \hline
5 & Үндсэн хамгаалалт & Нэгдсэн хуваарийн дагуу & \\ \hline
\end{tabularx}

\vspace{0.5cm}
\begin{center}\textbf{Ажлын товч дүгнэлт}\end{center}
\makebox[\textwidth]{\dotfill}\\[0.3cm]
\makebox[\textwidth]{\dotfill}\\[0.3cm]
\makebox[\textwidth]{\dotfill}\\[0.3cm]
\makebox[\textwidth]{\dotfill}\\[0.6cm]
Удирдагч: \makebox[5cm]{\dotfill} /\supname/

\begin{center}
\vspace{0.6cm}
\textbf{ЗӨВШӨӨРӨЛ}
\end{center}
Оюутан \shortname-ийн бичсэн төгсөлтийн ажлыг ШУТИС-д хамгаалуулахаар тодорхойлов.\\[0.4cm]
Тэнхимийн эрхлэгч: \makebox[5cm]{\dotfill} /\chairname/
\end{titlepage}

\begin{titlepage}
\small
\begin{center}
{\scshape \univname\par}
{\scshape \facname\par}\vspace{0.3cm}
\textbf{ШҮҮМЖЛЭГИЙН ХУУДАС}
\end{center}
Компьютерын ухааны тэнхимийн төгсөх курсийн оюутан \shortname-ийн «\ttitle» сэдэвт төгсөлтийн ажлын шүүмж.

\newcommand{\reviewlines}{\makebox[\textwidth]{\dotfill}\\[0.15cm]\makebox[\textwidth]{\dotfill}\\[0.3cm]}
\begin{enumerate}
\item Сэдвээр дэвшүүлсэн асуудал, үүнтэй холбоотой онолын материалыг яаж судалсан байдал.\\ \reviewlines
\item Ажлын ерөнхий агуулга. Шийдсэн зүйлүүд, гаргасан үр дүн.\\ \reviewlines
\item Эмх цэгцтэй, стандарт хангасан байдал.\\ \reviewlines
\item Давуу тал болон дутуу зүйлүүд.\\ \reviewlines
\item Анхаарч тэмдэглэх зүйлүүд.\\ \reviewlines
\item Ерөнхий оноо. (5 оноо)\\ \reviewlines
\end{enumerate}
Шүүмж бичсэн: \makebox[5cm]{\dotfill} /\readname/\\[0.2cm]
Ажлын газар: \makebox[9cm]{\dotfill}\\[0.2cm]
Хаяг (Утас): \makebox[9cm]{\dotfill}
\end{titlepage}

% Гарын үсэгтэй scan хуудсууд (байгаа бол %-ийг арилга):
%\includepdf[pages=-]{signatures/huudas1.pdf}
%\includepdf[pages=-]{signatures/huudas2.pdf}
%\includepdf[pages=-]{signatures/huudas3.pdf}
